\section{Introduction}
\label{sec:intro}

Large language models are now embedded in applications that read untrusted text,
call tools, and act on behalf of users. Two failure modes dominate the security
discussion. \emph{Jailbreaks} coax a model into producing content its operator
intended to withhold~\cite{zou2023universal,shen2023dan,chao2023pair,yi2024jailbreaksurvey}. \emph{Prompt
injection} is subtler: adversarial instructions embedded in data the model
processes—an email, a web page, a tool result—override the operator's intended task,
so that an output that violates the intended instruction hierarchy can nonetheless
look benign~\cite{greshake2023injection,liu2024formalizing,wallace2024instruction}.
Both threats have motivated a wave of defenses, from single-model hardening through
preference optimisation and structured
inputs~\cite{chen2025secalign,chen2024struq} to runtime guardrails and moderation
classifiers~\cite{inan2023llamaguard,rebedea2023nemo,markov2023holistic,zeng2024shieldgemma}.

A distinct and increasingly popular design filters LLM outputs with a \emph{team} of
LLM agents rather than a single model. AutoDefense~\cite{zeng2024autodefense} is the
canonical example: an Input agent forwards a candidate response to an Intention
Analyzer and an Original-Prompt Analyzer, a Judge agent renders a verdict, an optional
Llama-Guard agent~\cite{inan2023llamaguard} adds a content check, and a Coordinator
combines these into a final allow/block decision enforced by an Output agent. The
reported effect is substantial: on one configuration AutoDefense reduces the jailbreak
attack success rate from $55.74\%$ to $7.95\%$ using a three-agent LLaMA-2-13b
committee~\cite{zeng2024autodefense}. Multi-agent scaffolding of this kind—role
specialisation, analysis, and adjudication—mirrors the broader move toward LLM agent
societies~\cite{wu2023autogen,hong2023metagpt,li2023camel,qian2023chatdev,du2023debate,xi2023agentsurvey}.

\paragraph{The blind spot.}
These pipelines secure the content they inspect, but they do not secure
\emph{themselves}. The defense is a multi-agent system, and it inherits every threat
that multi-agent LLM systems are now known to
face~\cite{lee2024promptinfection,tian2023evilgeniuses,ju2024flooding,he2025redteaming,zhang2024psysafe,zhang2024asb}.
Concretely, an AutoDefense-style agency makes four assumptions that an adversary can
attack:
\begin{enumerate}[leftmargin=1.6em,itemsep=2pt,topsep=2pt]
  \item \textbf{Honest defenders.} Every internal agent is assumed faithful. A single
        compromised or fine-tuned-to-collude judge, or a supply-chain-tainted analyzer,
        is outside the model.
  \item \textbf{A trusted coordinator.} Verdicts flow through one aggregation point with
        no voting, quorum, or Byzantine tolerance; whoever controls it controls the
        decision.
  \item \textbf{Trustworthy self-inspection.} The untrusted candidate is fed directly to
        the judges. But the judges are themselves LLMs, and text crafted to look like a
        response can carry instructions aimed at \emph{them}—a \emph{second-order} prompt
        injection against the defenders, precisely the attack that flips
        LLM-as-a-judge verdicts~\cite{shi2024judgeinjection,zheng2023judging,wang2023notfair}.
  \item \textbf{Static adversaries.} Evaluation uses fixed attack sets, not an attacker
        that adapts to the pipeline or its aggregation rule—exactly the setting where
        robust aggregators have historically
        broken~\cite{fang2020poisoning,baruch2019little,xie2019fall}.
\end{enumerate}
A single hardened model such as SecAlign~\cite{chen2025secalign} has no co-defender to
collude with and no coordinator to subvert; these threats appear only once the defense
is a team. The gap is therefore specific to multi-agent defenses and cannot be closed
by swapping in a better single model.

\paragraph{This paper.}
We treat the defense pipeline as a first-class attack surface and ask what it takes to
tolerate a compromised minority of its own agents while resisting injection from the
content it reads. Our answer, \Aegis, has four components, illustrated in
\cref{fig:architecture}:
\begin{enumerate}[leftmargin=1.6em,itemsep=2pt,topsep=2pt]
  \item \emph{Byzantine-robust verdict aggregation.} The single coordinator is replaced
        by a robust aggregation rule over judge verdicts—Krum-style
        selection~\cite{blanchard2017krum}, coordinate-wise
        median~\cite{yin2018byzantine}, or geometric
        median~\cite{chen2017byzantine,pillutla2022rfa}—so that a malicious minority
        cannot flip the decision.
  \item \emph{Payload isolation.} The untrusted candidate is confined to a structured
        data channel, separated from each judge's instruction channel following the
        structured-query and spotlighting
        principles~\cite{chen2024struq,hines2024spotlighting}, closing the second-order
        injection surface the defense creates by reading untrusted text.
  \item \emph{Hardened judges.} Each judge is individually fine-tuned to resist
        injection~\cite{chen2025secalign,chen2024struq}, so that isolation and
        aggregation defend a committee whose members are already robust.
  \item \emph{Task-integrity generalisation.} The verdict target is lifted from
        ``harmful content'' to ``violated the intended task or privilege
        hierarchy''~\cite{wallace2024instruction}, extending the pipeline from jailbreak
        filtering to prompt-injection filtering.
\end{enumerate}

\paragraph{What we claim, and what we do not.}
We are deliberately conservative. Robust aggregation borrowed from distributed learning
comes with a tolerance condition ($2f{+}2<n$ for
selection~\cite{blanchard2017krum}, $f<n/2$ for median
rules~\cite{yin2018byzantine}), and its guarantees are conditional on the honest
verdicts actually concentrating—an assumption that is vacuous if every judge shares a
backbone and fails identically~\cite{kuncheva2003diversity}. We therefore prove a
decision-integrity guarantee \emph{under explicit assumptions}, and separately
characterise how it decays as judge failures become correlated. We do not claim that a
lower attack success rate under compromise equals security; we claim \emph{graceful
degradation} and we report residual failures. Our theoretical results are stated with
their assumptions in full, and the empirical numbers in this manuscript are placeholders
that fix the evaluation protocol rather than measured outcomes.

\paragraph{Contributions.}
\begin{itemize}[leftmargin=1.4em,itemsep=2pt,topsep=2pt]
  \item \textbf{Threat model (\cref{sec:system,sec:problem}).} The first threat model,
        to our knowledge within the verified literature, that drops AutoDefense's
        honest-agent and single-coordinator assumptions and adds second-order injection
        against the judges, with attacker knowledge, capability, collusion, and adaptivity
        made explicit.
  \item \textbf{Architecture (\cref{sec:method,sec:algorithm}).} \Aegis, combining
        Byzantine-robust verdict aggregation, payload isolation, and hardened judges,
        generalised from jailbreak to task-integrity, together with the aggregation
        algorithm over the semantic verdict space.
  \item \textbf{Theory (\cref{sec:theory,sec:complexity}).} A decision-integrity theorem
        for robust verdict aggregation under honest-concentration and margin assumptions,
        a correlated-failure analysis that makes the diversity requirement quantitative,
        and a security-versus-cost complexity analysis.
  \item \textbf{Adaptive evaluation (\cref{sec:experiments}).} An evaluation protocol
        against compromised judges, colluding minorities, second-order payload injection,
        and adaptive attacks on the aggregation rule—absent from prior multi-agent
        defenses—with baselines AutoDefense, a single hardened model, and plain majority
        vote, and a stated correlated-failure test over homogeneous versus diverse
        backbones.
\end{itemize}

\Cref{sec:related} positions the work; \cref{sec:prelim} fixes preliminaries;
\cref{sec:system,sec:problem} give the system and threat model and problem formulation;
\cref{sec:method,sec:algorithm} present the architecture and algorithm;
\cref{sec:theory,sec:complexity} develop the theory and cost analysis;
\cref{sec:experiments} specifies the evaluation; \cref{sec:discussion,sec:limitations}
discuss implications and limitations; \cref{sec:conclusion} concludes.
